%% paper.tex -- LaTeX conversion of paper.md (v0.8), format conversion only.
%% Source: github.com/Oliverds321/RH-72. pdflatex-compatible; pure ASCII.
\documentclass[11pt]{amsart}
\usepackage[T1]{fontenc}
\usepackage{lmodern}
\usepackage{amsmath}
\usepackage{amssymb}
\usepackage{amsthm}
\usepackage{url}
\urlstyle{same}
\makeatletter
\g@addto@macro{\UrlBreaks}{\UrlOrds}
\makeatother

\allowdisplaybreaks
\emergencystretch=2em

%% Operators
\newcommand{\re}{\operatorname{Re}}
\newcommand{\im}{\operatorname{Im}}
\newcommand{\tr}{\operatorname{tr}}
\newcommand{\supp}{\operatorname{supp}}
\newcommand{\conj}{\operatorname{conj}}
\DeclareMathOperator{\frobSq}{frobSq}

%% Theorem environments: the paper's own numbering is preserved exactly
%% (statements are declared with their fixed numbers; nothing is auto-numbered).
\theoremstyle{plain}
\newtheorem*{theoremone}{Theorem 1}
\newtheorem*{corollarytwo}{Corollary 2}
\newtheorem*{corollarythree}{Corollary 3}
\newtheorem*{propositionthreeone}{Proposition 3.1}
\newtheorem*{lemmafourone}{Lemma 4.1}
\newtheorem*{lemmafourtwo}{Lemma 4.2}
\newtheorem*{lemmafourthree}{Lemma 4.3}
\newtheorem*{lemmafourfour}{Lemma 4.4}
\newtheorem*{lemmafourfive}{Lemma 4.5}
\newtheorem*{lemmafiveone}{Lemma 5.1}
\newtheorem*{lemmafivetwo}{Lemma 5.2}
\newtheorem*{lemmafivethree}{Lemma 5.3}
\newtheorem*{lemmasfivetwoprime}{Lemmas 5.2$'$/5.3$'$}
\newtheorem*{lemmasixone}{Lemma 6.1}
\newtheorem*{lemmaeightone}{Lemma 8.1}
\newtheorem*{lemmaeighttwo}{Lemma 8.2}
\theoremstyle{remark}
\newtheorem*{remarkenv}{Remark}

\begin{document}

\title[Simple zeros for the primitive Dirichlet family]{Simple
zeros on the critical line for the primitive Dirichlet family: 72.128\% of the
Generalized Riemann Hypothesis, on average, unconditionally}
\author{Oliver D'Souza}
\date{2026-08-18}
\thanks{Preprint --- 2026-08-18 (v0.8).}
\thanks{Repository: \url{https://github.com/Oliverds321/RH-72} (scripts, logs, data).}

\begin{abstract}
We prove that, unconditionally, at least 0.72128 of the zeros of Dirichlet
L-functions $L(s, \chi)$, averaged over all primitive characters of modulus $q \le Q$
in the ordinate window $[T, 2T]$ with $T = (\log Q)^{3+\varepsilon}$, are simple and lie
on the critical line, with an explicit rate of convergence. On the corresponding
families this improves the unconditional records --- 56\% over all primitive characters
at $q \asymp Q$ (Conrey--Iwaniec--Soundararajan [CIS2], at their window floor
$T \ge (\log Q)^{6}$, which our statement reaches at $r \ge 6$, and up to window shape:
[CIS2] counts $|\gamma| \le T$ where we count dyadic blocks) and 0.6044 over the even
primitive characters (Sono [So21]) --- by 14.99 and 8.75 points; both partners carry a
smooth conductor weight where our count is sharp and unweighted. Per character, the
strongest PUBLISHED unconditional proportion is Wu's 0.4074 [Wu19]; the 0.672500703679
of the August 2026 preprint [R], on which our architecture builds and which is at this
writing unposted to any preprint server and unrefereed, is likewise per character, in a
range of $T$ relative to $q$ disjoint from ours. The $q \le Q$ family of Theorem 1, in
which every primitive character of every modulus up to $Q$ is counted with equal weight,
has no counterpart in the literature. Our results consume six statements imported from
[R], isolated in \S2.1, each a sorry-free theorem in [R]'s public Lean 4 artifact ---
at [R]'s own parameter class; their use at THIS paper's parameters (bandwidth
$\lambda > 1$) rests on the file-level audit of \S9, and the re-parameterisation that
would certify it mechanically is an open obligation of the formalization track. The
input is one-sided: the multiplicative large sieve --- lossless precisely when the
family is larger than the polynomial, which holds here at every bandwidth below 2 ---
replaces the GRH pair-correlation asymptotic at the cost of an explicit constant, and a
$C$-penalised variational problem converts that constant into a proportion. No
asymptotic large sieve and no prime-correlation asymptotic is used, and no unproved
conjecture is assumed: beyond classical prime number theory, every input is one of the
six imported statements of \S2.1 --- proved theorems of [R], consumed here as premises.
The per-character architecture (the rank--trace certificate and the Weil-form grid
Gram) is imported from the August 2026 preprint [R]; everything family-level is new.
\end{abstract}

\maketitle

\section{Introduction}

\subsection{The result}

Let $\mathfrak{F}_Q$ denote the set of primitive Dirichlet characters of modulus
$q \le Q$, and for $\chi \in \mathfrak{F}_Q$ let $N_\chi(T, 2T)$ count the nontrivial
zeros $\rho = \beta + i\gamma$ of $L(s, \chi)$ with $T < \gamma \le 2T$, and
$N^{s}_{0,\chi}(T, 2T)$ those that are simple and on the critical line. The
architecture we build on does not assume the zeros lie on the line, nor that $\gamma$
is real for $\beta = 1/2$ --- handling possible off-line zeros unconditionally is the
point.

\begin{theoremone}
Fix $\varepsilon > 0$. There are $Q_0(\varepsilon)$ and an effectively computable
$c(\varepsilon)$ such that for $Q \ge Q_0(\varepsilon)$ and
$T = T(Q) = (\log Q)^{3+\varepsilon}$, unconditionally,
\[
\sum_{\chi\in\mathfrak{F}_Q} N^{s}_{0,\chi}(T, 2T)
\;\ge\; \bigl(P - c(\varepsilon)\cdot(\log \log Q)/(\log Q)\bigr)
\cdot \sum_{\chi\in\mathfrak{F}_Q} N_\chi(T, 2T),
\]
where {\boldmath$P = 0.7212835668\dots$} is the value of an explicit variational
problem (\S10) at $C = \pi^4/18$. (The same statement holds at
$T = (\log Q)^{r+\varepsilon}$ for every $r \ge 3$, with $c = c(r, \varepsilon)$
effectively computable; a computed --- not proved --- value of the construction's
budget constant is in the remark of \S10.3, outside this statement.)
\end{theoremone}

\begin{corollarytwo}[dyadic]
The same holds for the dyadic subfamily $q \in (Q/2, Q]$ with
{\boldmath$P_{\mathrm{dyad}} = 0.7099167448\dots$} ($C = 2\pi^4/27$,
$\lambda^* = 1.1931581210$), for every $T = (\log Q)^{r+\varepsilon}$, $r \ge 3$. The
comparison partner in the literature is Conrey--Iwaniec--Soundararajan [CIS2] (all
primitive characters, $q \asymp Q$, no parity restriction, weighted $\Psi(q/Q)/\varphi(q)$
--- the weight largely cancels in a ratio of two identically weighted sums, though we do
not claim it exactly does): their theorem gives $14/25 = 56\%$ simple-and-on-line
(58.65\% announced at optimised parameters), in the window $|\gamma| \le T$ with floor
$T \ge (\log Q)^{6}$. At $r \ge 6$ our statement enters [CIS2]'s window class
(asymptotic budget constant $\approx 4.6$ rather than $\approx 3.3$; at finite $Q$ the
$r = 6$ design is in fact STRONGER, since the buffer and ramp rows fall faster in $T$
than the zone row rises --- \S10.4), and there the comparison is like-for-like up to
window shape ([CIS2] counts all $|\gamma| \le T$; we count the dyadic block $[T, 2T]$
--- the two statements are never literally identical, since this method needs
$T \ge (\log Q)^{3}$ and cannot cover the low blocks): \textbf{+14.99 points} over
$14/25$ (+12.34 over the announced optimum). Theorem 1's $q \le Q$ family has no
comparison partner in the literature.
\end{corollarytwo}

\begin{corollarythree}[even primitive, dyadic --- the Sono-comparable statement]
Sono's family [So21] is the EVEN primitive characters of dyadic conductor. The even
subfamily carries exactly half the character mass ---
$|\mathfrak{F}_{\mathrm{even}}| = \tfrac{1}{2}|\mathfrak{F}| + O(Q)$, provably:
$\#\text{even-prim}(q) = (\varphi^{*}(q) + S(q))/2$ with
$S(q) := \sum_{\chi \text{ prim mod } q} \chi(-1)$ multiplicative ($S(p) = -1$ for odd
$p$, $S(p^a) = 0$ for odd $p$ and $a \ge 2$; $S(2) = 0$, $S(4) = -1$, $S(2^a) = 0$ for
$a \ge 3$), so $S$ is supported on the squarefree-odd $q$ and
$4\cdot(\text{squarefree-odd})$ and $\sum_{q\le Q} S(q) = O(Q)$. The subfamily passage
is a TWO-ZONE argument, and the two zones use different mechanisms --- positivity
alone, applied in both zones, would double the in-zone coefficient as well and prove
only 0.3693 (\S11 at in-zone weight $2|\alpha|$), which is NOT this corollary. IN-ZONE
($n, m \le Y = Q^{1-\delta'}$), the even-subfamily main term is computed exactly by
parity projection, $\sum_{\chi \text{ even}} = \tfrac{1}{2}\sum^{*}_{\chi}(1 + \chi(-1))$:
the 1-half is half the full-family form, in step with the halved denominator, and the
$\chi(-1)$-half pairs $n \equiv -m \pmod q$ and is governed by the aggregated identity
in its $n + m$ form, Lemma 5.2$'$ (stated in \S5) --- in-zone
$n + m \le 2Q^{1-\delta'} \ll Q$, so Lemma 5.3$'$'s divisor bound applies and it is
negligible at the same zone-edge calibration. The in-zone coefficient of the
variational problem is therefore UNCHANGED. OUT-ZONE (the $C$-penalised region), where
$n$ reaches $X = Q^\lambda$ and $n + m$ outruns every modulus --- exactly where the
projection route fails by $Q^{\lambda-1}$ (\S12.3) --- the sieve input passes to the
subfamily by positivity of its summands: the budget $N + Q^2 - 1$ is unchanged, the
denominator halves, and the OUT-ZONE constant doubles, $C \to 2C$. The error rows
(ends, cross, tail) are one-sided or absolute-valued, so the subfamily is bounded by
the full family at the cost of a factor $\le 2$ against the halved denominator --- each
is priced with an order-of-magnitude margin in the budget, and doubling them moves
nothing. Finally the denominator is a ZERO count, not the character count just proved:
$\mathcal{N}_{\mathrm{even}} = \tfrac{1}{2}\mathcal{N} + O(QT\mathcal{L})$, by partial
summation of
$N_\chi(T, 2T) = (T/2\pi)(\log(qT/2\pi) + 2 \log 2 - 1) + O(\log qT)$ against
$\sum_{q\le x}S(q) = O(x)$ --- relative error $O(1/Q)$. (Parity uniformity holds where
it is touched at all: $\mu_q$'s parity dependence is exactly
$\pi \re \cot(\pi(\tfrac{1}{4} + i\tau/2)) = O(e^{-\pi\tau})$ by the reflection formula
--- beyond floating point at $\tau \asymp T$.) The argument with out-zone constant $2C$
gives, for even primitive $\chi$ with $q \in (Q/2, Q]$:
{\boldmath$P_{\mathrm{even,dyad}} = 0.6919434301\dots$} ($C = 4\pi^4/27$,
$\lambda^* = 1.1015998422$) --- \textbf{+8.75 points over the unconditional family
record 0.6044 [So21]}. (Even primitive over $q \le Q$: 0.6980745436 at $C = \pi^4/9$.)
The identical argument with the projector $\tfrac{1}{2}\sum^{*}(1 - \chi(-1))$ gives
the ODD primitive family at the same constants:
{\boldmath$P_{\mathrm{odd,dyad}} = 0.6919434301$} ($q \in (Q/2, Q]$; 0.6980745436 over
$q \le Q$) --- a class in which the literature contains no result at all ([So21] is
stated for even characters only, in its (1.8)). Remark: a parity-refined large sieve at
budget $\tfrac{1}{2}(N + Q^2 - 1)$ would restore $C$ in the out-zone, but the
$n \equiv -m$ orthogonality route --- which IS what the corollary uses in-zone ---
fails OUT-ZONE at $\lambda > 1$ by $Q^{\lambda-1}$, the same mechanism as the cross
term of \S4, and no parity-restricted sieve with the $\tfrac{1}{2}$ constant appears to
be available; mirroring Sono's smooth weight $W \le 1$ on $[1, 2]$ would further
inflate $C$ by $3/(2\int W(x)x \, dx)$.
\end{corollarythree}

Both statements are family averages: they say nothing about any individual $L(s, \chi)$,
and a zero-density subfamily could in principle carry all its zeros off the line. This
is the same statement class as [CIS2] and [So21] ([So16] is GRH-conditional and counts
low-lying zeros weighted by $|\widehat{\Phi}(i\gamma)|^2$ for a fixed $\Phi$, so its
count concentrates at $|\gamma| \ll 1/\log Q$ --- a different class on both grounds;
the harmonic $1/\varphi(q)$ weight is NOT a class separator, since [CIS2] itself
carries $\Psi(q/Q)/\varphi(q)$ --- see Corollary 2's statement).

\subsection{What was known}

Per character, the strongest PUBLISHED unconditional proportion of simple-and-on-line
zeros of Dirichlet L-functions is Wu's \textbf{0.4074} [Wu19] (with 0.4172 on the
line), for every primitive $\chi \bmod q$ with $\log q = o(\log T)$; for the family of
all primitive characters of a single FIXED modulus, Dickinson [Di24] gives 38.2\% on
the line, at $T \gg q^\varepsilon$; for $\zeta$ alone the corresponding constants are
Pratt--Robles--Zaharescu--Zeindler's 0.407511 and 0.417293 [PRZZ] --- the latter is the
``more than five-twelfths'' quoted as 41.7\% in the recent expository literature
[GS26]. Higher still, but not published: the preprint [R], dated 11 August 2026 and
revised in place 13 August 2026, proves for $\zeta$ and --- its \textbf{Theorem B} ---
for every FIXED primitive $\chi$ that at least
$3/2 - (1/\sqrt{2})\cot(1/\sqrt{2}) = \textbf{0.672500703679}\dots$ of the zeros in
dyadic windows $[T, 2T]$ are simple and on the critical line, at full bandwidth
$\lambda = 1$ with the Montgomery--Taylor window and $T \ge T_0$ depending on the
character ([R]'s Remark 6.1: no $\lambda < 1$ improves the constants). It is
Lean-formalized (the artifact retains the lettering of the 10 August version, in which
the Dirichlet case is Theorem E); it is not posted to a preprint server, and we are
aware of no journal submission, refereeing, erratum or withdrawal. The constant itself
is not new: it appears in [BGST2, Thm 2] under a narrow-vertical-box hypothesis, and
[R]'s contribution is removing that hypothesis. We make no claim about [R]'s reception.
\textbf{Our construction is not a corollary of [R]'s Theorem B}: that theorem is per
character, with $T_0$ and the error constants depending on $q$ --- [R] states the
Dirichlet case with $O_q(\cdot)$ errors and claims no uniformity --- and it supplies no
family statement at $T = \operatorname{polylog}(Q)$; the uniformity in $q$ is the work,
and it is where every difficulty of this paper lives. The same disjointness applies to
[Wu19], whose hypothesis $\log q = o(\log T)$ is the exact reverse of our range, in
which $\log T \asymp r \log \log Q = o(\log q)$ for almost every member of the family.
We record for orientation that our family average 0.72128 exceeds both per-character
constants (by 4.88 and 31.39 points), but the statements are of different kinds --- a
per-character bound implies the family average in its own range, not conversely --- and
we make no claim about any individual $L(s, \chi)$.

On family average the unconditional records are, by class: over the EVEN primitive
characters of dyadic conductor, Sono's 0.6044 simple-and-on-line (0.6107 on-line)
[So21]; over ALL primitive characters at $q \asymp Q$ with the harmonic weight
$\Psi(q/Q)/\varphi(q)$, Conrey--Iwaniec--Soundararajan's $14/25 = 56\%$ (58.65\%
announced at numerically optimised parameters) [CIS2]; and, for the weaker quantity
``simple, critical line not required'' over that same weighted family, Wu's 0.60261,
or 0.66433 under GRH [Wu16]. Under GRH the low-lying benchmark is [So16]'s 93.22\% ---
a different quantity again (\S1.1). On the corresponding families, Corollary 2 exceeds
[CIS2] by \textbf{+14.99 points} (at $r \ge 6$, inside [CIS2]'s $(\log Q)^{6}$ window
floor and up to window shape and the harmonic weight --- see the corollary) and
Corollary 3 exceeds [So21] by \textbf{+8.75 points}. Since simple-and-on-line implies
simple, Corollary 2 also gives $\ge 0.70992$ simple, exceeding [Wu16]'s GRH figure
unconditionally by +4.56 points, modulo his harmonic weight. (Each corollary pays its
own family's constant; \S12.2--12.3.) On the procedural point: dependence on unrefereed
work is not unusual in this literature --- the asymptotic large sieve [CIS1] is listed
by its first author among unpublished papers and has remained so since 2011, and
[CIS2], [CLLR], [So16] (via [CLLR]), [CL] and [Wu16] all rest on it, [Wu16] citing it
by its arXiv number as unpublished. \S2.1 isolates every statement this paper imports
from [R], so the dependency is checkable.

\subsection{The idea, and the one-sided large sieve}

The certificate of [R] consumes, for each $\chi$, a second moment of a zero-counting
linear form --- a family pair-correlation-type quantity --- \textbf{with a single sign:
an upper bound suffices.} Under GRH the relevant object $F(\alpha)$ is known
asymptotically for the Dirichlet family: $F(\alpha) = \min(|\alpha|, 1)$ on
$|\alpha| < 2$ --- for the PRIMITIVE family this is Chandee--Lee--Liu--Radziwi\l\l{}
[CLLR, Thm 2] (their Corollary 3 is the character-sum consequence); the antecedent,
over ALL characters rather than the primitive ones, is \"Ozl\"uk [\"Oz] --- the
primitivity is exactly the defect [CLLR] name and fix. Our Corollary 2 produces a
one-sided unconditional counterpart of that statement's CONTENT for the dyadic family,
where the normalisation matches exactly --- for the $q \le Q$ family of Theorem 1 the
consumed object is the large-sieve majorant of \S12.1, not [CLLR]'s $F_\Phi$.
Unconditionally, in the $t$-aspect, the trivial bound loses a factor $\log T$
([BGST, Thm 1] --- the unconditional $t$-aspect Montgomery theorem --- and [Go],
survey). The observation this paper runs on --- we are not aware of it being recorded,
and we state it as a mechanism rather than a theorem --- is that \textbf{on family
average the loss is a constant}: the multiplicative large sieve
\[
\sum_{q\le Q} \sum^{*}_{\chi \bmod q} \Bigl|\sum_{n\le N} a_n \chi(n)\Bigr|^2
\le (N + Q^2 - 1)\cdot\sum|a_n|^2
\]
is lossless precisely when the polynomial is shorter than the family, $N \ll Q^2$ ---
which holds in the $q$-aspect for every bandwidth $\lambda < 2$ and never in the
$t$-aspect. The price of one-sidedness is the constant
$C = Q^2/|\mathfrak{F}_Q| = \pi^4/18 + o(1)$ --- the ``smaller than $Q^2$ by a
constant factor'' of [CIS1] --- paid on the off-diagonal range $1 < |\alpha| \le \lambda$
where GRH-truth would be $F = 1$. The certificate turns that overshoot into a
proportion via a $C$-penalised variational problem (\S10) whose value exceeds the
record's constant for every finite $C$.

\begin{remarkenv}[the size of the margin]
Corollary 2 exceeds [CIS2] by fifteen points using a sieve inequality available since
1974, and the margin is large for this literature. The structural explanation is
short: the certificate the sieve plugs into is days old, so the combination could not
have been tried before. We do not ask the reader to weigh explanations --- a result of
this profile can be wrong at any margin, and the repository is built so that checking,
not argument, settles it.
\end{remarkenv}

\subsection{Why $q \le Q$, against the field's convention}

Every family result in this literature normalises $q \asymp Q$ (we verified fifteen
papers without exception --- [\"Oz], [CIS1], [CIS2], [CLLR], [So16], [So21], [Wu16],
[CL], [DPR], [FM], [Mi], [ILS], [GZ], [BP], [LM]; the restriction lives in the smooth
weight's support, and [CIS1] states the convention verbatim: ``the conductor $q$ is
restricted by a smooth function $\Psi$ of compact support in $\mathbb{R}_{+}$, so
$q \asymp Q$'' --- [CIS2] and [Wu16] inherit $\Psi$ from it). The convention is
inherited from machinery: the asymptotic large sieve [CIS1] --- the engine under all
the asymptotic results --- is stated for $q \asymp Q$. \textbf{This paper does not use
the asymptotic large sieve.} (The polylog-$T$ window, by contrast, is NOT part of the
novelty claim: [So21] already operates from the weaker floor $(\log Q)^{2}$, below our
$(\log Q)^{3+\varepsilon}$. What no prior result covers is the FAMILY --- every
primitive character of every modulus $q \le Q$, at equal weight.) The classical
sieve's native and sharp range is exactly $q \le Q$; restricting to the dyadic family
costs (the budget $N + Q^2 - 1$ is charged against the largest modulus while the
dyadic family supplies only $3/4$ of the characters), which is precisely why
$C_{\mathrm{dyad}} = 2\pi^4/27 > \pi^4/18$. Our family choice is a design argument,
not a liability. The price is a normalisation discipline (\S12): the consumed object
is a large-sieve majorant of a family second moment at the common scale
$\mathcal{L} = \log(QT/2\pi)$ --- it is NOT the form factor $F_\Phi$ of [CLLR], which
is correctly normalised only at $q \asymp Q$; the conductor-average
$\langle \log q\rangle_{\varphi^{*}} = \log Q - 1/2 + O(1/Q^{\dots})$ is stated with
its constant, in the tradition of Fiorilli--Miller [FM] and Miller [Mi, (2.8) and
Rem. 3.2], who compute resp.\ recommend exactly this substitution.

\subsection{Honesty about effectivity}

Theorem 1 is an asymptotic statement, as is every prior result in this line; unusually,
we have computed our construction's finite-$Q$ content, and we state it exactly once,
in a remark labelled ``for orientation'' (\S10.4). Nothing effective is claimed
anywhere else in the paper, and no effective figure appears in this introduction or in
the abstract by design.

\subsection{Ancillary computations}

The numerical constants in this paper (the variational values of \S10--11 and the
orientation figures of \S10.4) were computed by scripts included as ancillary files,
whose identities and conventions are cross-checked against toy-scale exact
computations and a dataset of 5,230 zeros of the 60 primitive L-functions of moduli
5--19; the variational constants are reproduced to ten digits by three independent
implementations (the high-precision closed-form/Nystr\"om pair is shipped as
\texttt{payoff\_hp2.py}; the coarse grid anchor \texttt{qg\_check.py} reaches 4--5
digits, and its printed $v \ge 0$ flag is a discretisation artifact of its own
edge-sampling scheme: on an exact-cell discretisation $\min v$ is POSITIVE at every
tested resolution and decays like the grid spacing --- $v$ vanishing linearly at the
free boundary --- with the ladder $n = 500\dots8000$ shipped in
\texttt{ext\_checks.py}/\texttt{log\_ext\_checks.txt}; the closed-form solve is
positive with strict KKT complementarity). The orientation figures of \S10.4 are
computed by the shipped \texttt{table2\_check.py}. These checks validate transcription
and conventions; they are not part of any proof.

\section{Notation, standing conventions, and imported hypotheses}

\subsection{What is imported from [R], as numbered inputs}

Everything this paper uses from the unrefereed preprint [R] is one of the following
six statements, each a sorry-free theorem in [R]'s Lean artifact (file references in
the ancillary notes): \textbf{(H1)} the rank--trace certificate: for the direct sum of
Hermitian blocks, $\#\{\text{simple, on-line}\} \ge 4 \tr \widehat{G} -
\lVert\widehat{G}\rVert^2_F - 2\mathcal{N} - 3N_{\mathrm{II}} -
(\text{perturbation loss})$, linear in $(\tr, \lVert\cdot\rVert^2_F)$; \textbf{(H2)}
the Weil-form grid Gram identity $Gz = Gp$ per primitive $\chi$ with density
$\nu_{X,\chi} = \mu_q + P_{X,\chi}$ and no pole term; \textbf{(H3)} the pointwise
positivity $g = \varphi^2\star\varphi^2 \ge 0$; \textbf{(H4)} the Gevrey-2 profile
$\varrho_2$ with derivative constants $(A, B) = (36/e, 2e^{8})$ and the ramp lemma
$\lVert\varphi^{(k)}\rVert_1 \le 2Bw(A/w)^k k^{2k}$; \textbf{(H5)} the $\nu$-generic
ends majorants (the cap-free L-intrinsic forms) with their explicit constants;
\textbf{(H6)} the $q$-uniform local count $N_\chi(t, t+1] \le A_0\cdot\log(q(|t|+3))$
with absolute $A_0$. The $\lambda \le 1$ audit: [R]'s $T$-aspect development
frequently assumes $\lambda \le 1$ (equivalently $X \le T$, which fails here at EVERY
$\lambda$ since $X = (QT/2\pi)^\lambda \gg T$); the audit of \S9 (file-level record in
the companion working notes) shows that none of H1--H6's proofs consumes it --- H1 is
bandwidth-free by its signature, H3--H6 are cap-free by construction, and the
cap-gated material ([R]'s prime side, ends instantiation, payoff algebra,
\texttt{calE} analytics) is replaced, not imported --- with one Lean-structural caveat
stated in \S9: [R]'s \texttt{Params.Valid} bundles $\lambda \le 1$ (and $w \ge 1$), so
formalization-track reuse requires the \texttt{ParamsQ} weakening first. This is a
LIVE obligation, named as a risk: the artifact certifies H1--H6 at [R]'s parameter
class ($\lambda \le 1$, $w \ge 1$, $D_0 = \sqrt{T}$); at this paper's parameters their
validity rests on the \S9 audit --- ``no proof consumes the gated fields'' --- until
the \texttt{ParamsQ} re-parameterisation compiles. The $w \ge 1$ clamp in \S10.3's
design is a case of the artifact's validity class already shaping the design. One
further artifact fact, for completeness: none of H1--H6 routes through [R]'s
\texttt{PairCeiling/} modules --- the one place the artifact carries an extra-Lean
hypothesis (\texttt{EnclOK}: 256 interval-arithmetic enclosures computed outside
Lean); \texttt{PairCeiling} is imported only by the artifact's top-level umbrella and
its axiom comparator, not by any module H1--H6 lives in (verified by import grep).

\subsection{Notation}

Primitive $\chi \bmod q$, $q \le Q$; $\kappa = \kappa(\chi) \in \{0, 1\}$ the parity.
$\mathcal{L} := \log(QT/2\pi)$; bandwidth parameter $\lambda \in (0, 2)$;
$L := \lambda\mathcal{L}$; $X := e^L = (QT/2\pi)^\lambda$. Window $I := [T, 2T]$,
$T = (\log Q)^{3+\varepsilon}$. Grid $\tau_k := T + kh$, $h := 2\pi/L$,
$d := \lfloor LT/2\pi\rfloor$ ($\asymp \mathcal{L}^{1+r+\varepsilon}$ ---
$\mathcal{L}^{4.5}$ at the design of record: polylog). paperFT (the Fourier normalisation of [R],
carried unchanged):
$h_f(z) := \int f(u)e^{izu}du$. Taper: the Gevrey-2 profile $\varrho_2$ of [R] (their
P1--P4: normalized primitive of $e^{-1/x}e^{-1/(1-x)}$, derivative constants
$(A, B) = (36/e, 2e^{8})$), $\varphi(u) := \varrho_2((L/2 - |u|)/w)$, ramp width $w$
and buffer $D_0$ --- the values $w = 3 \log \mathcal{L}$,
$D_0 = \mathcal{L}^2 \log \mathcal{L}$ displayed in \S7 are ONE admissible choice (the
REFERENCE design), at which every closing margin is proven with
$(\log \mathcal{L})^2$-room; the design of record for the budget is \S10.3's
jointly-optimised BALANCED design ($w = \Theta(1)$, $D_0$ balanced), whose closing
condition is met with equality by construction at the corner design --- \S7.2; NOT a
fortiori from the reference design's: the balanced design trades all its margin for
budget. $\Phi := h_{\varphi^2}$; $g := \varphi^2\star\varphi^2 \ge 0$;
$I' := (T - D_0, 2T + D_0]$. Zone edge
{\boldmath$\delta' := 3 \log \log Q/\log Q$}. Densities:
$\mu_q(\tau) := (1/2\pi)[\log(q/\pi) + \re \psi(\tfrac{1}{4} + \kappa/2 + i\tau/2)]$;
$P_{X,\chi}(\tau) := -(1/\pi) \re \sum_{n\le X} \Lambda(n)\chi(n)n^{-1/2-i\tau}$;
$\nu_{X,\chi} := \mu_q + P_{X,\chi}$ --- \textbf{no pole term} ($L(s,\chi)$ entire for
primitive $\chi \ne \chi_0$). $\mathcal{N} := \sum_\chi N_\chi(T, 2T) \asymp
Q^2T\mathcal{L}$. $D_T(v) := \int_I e^{i\tau v}d\tau$.
$|\mathfrak{F}_Q| = (18/\pi^4)Q^2(1 + o(1))$; $C := Q^2/|\mathfrak{F}_Q| \to \pi^4/18$
(dyadic: $2\pi^4/27$).

All parameter SHAPES are forced (the values within the admissible region are the
design's to choose --- \S10.3): $T \ge (\log Q)^{3}$ by the buffer/rate trade
(\S10.3); $D_0 \gg L$ by the tail (\S7); $w$ by the Gevrey closing condition;
$\delta'$ by the zone-edge calibration (\S5); the common taper (one coefficient vector
for the whole family) by the large sieve itself.

\section{The certificate}

Per $\chi$, the grid Gram $G(\chi)_{kl} := \int \widehat{\varphi}(\tau - \tau_k)
\widehat{\varphi}(\tau - \tau_l)\, \nu_{X,\chi}(\tau)d\tau$; by the $q$-uniform
explicit formula (\S9) this equals the zero-side matrix
$\sum_\rho m_\rho h_\varphi(\gamma_\rho - \tau_k)\cdot\conj(h_\varphi(\dots))$; hat
units $\widehat{G} := G/(aL^2)$, $a := L^{-1}\int\varphi^2$. The certificate,
displayed in full (H1; the $-2\mathcal{N}$ and $-3N_{\mathrm{II}}$ window terms are
part of the inequality, not corrections):
\begin{multline*}
\sum_\chi N^{s}_{0,\chi}(T, 2T) \;\ge\; 4 \tr \widehat{G}_{\mathrm{fam}}
- \lVert\widehat{G}_{\mathrm{fam}}\rVert^2_F - 2\mathcal{N}
- 3\cdot N_{\mathrm{II,fam}} \\
- \bigl[4B_{\mathrm{tr}} + 2B_F\cdot\lVert\widehat{G}_{\mathrm{fam}}\rVert_F
+ B_F^2\bigr],
\end{multline*}
linear in $(\tr, \lVert\cdot\rVert^2_F)$, so the direct sum over $\mathfrak{F}_Q$
aggregates identically --- with exactly one exception, the tail perturbation, which
aggregates by the pair mechanism of \S7.3.

\begin{propositionthreeone}[quantitative pair moment certificate]
Suppose for a given $(Q, T)$: (i) the display above; (ii)
$\mathcal{N}(1 - r_1) \le \tr \widehat{G}_{\mathrm{fam}}$ and
$\lVert\widehat{G}_{\mathrm{fam}}\rVert^2_F \le (\kappa_C + r_2)\mathcal{N}$, where
$\kappa_C$ is the Frobenius main-term constant of \S\S4--6 (NOT the parity
$\kappa(\chi)$ of \S2.2); (iii) $N_{\mathrm{II,fam}} \le r_3\mathcal{N}$; (iv)
$B_{\mathrm{tr}} \le r_4\mathcal{N}$ and $B_F \le r_5\sqrt{\mathcal{N}}$. Then
\[
\sum_\chi N^{s}_{0,\chi} \;\ge\; \bigl[2 - \kappa_C - (4r_1 + r_2 + 3r_3 + 4r_4
+ 2r_5\sqrt{\kappa_C+r_2} + r_5^2)\bigr]\cdot\mathcal{N},
\]
so the rate of Theorem 1 is carried by the explicit budget of \S10.3 (whose rows are
exactly the $r_i$), not lost to an $\varepsilon$-limit.
\end{propositionthreeone}

\emph{Proof:} the moment algebra of [R]'s certificate with the scalar $B$ split into
$(B_{\mathrm{tr}}, B_F)$ --- trace linearity, the Frobenius triangle inequality, and
$\sqrt{\frobSq(\widehat{E})} \le B_F$; three displays. $\blacksquare$ (The moment
inputs are \S\S4--9.)

The prime side must deliver: the trace $\approx \mathcal{N}$ (\S9), and
$\frobSq \le (\kappa_C + o(1))\cdot\mathcal{N}$ with $\kappa_C =$ the bracket the
variational problem penalises (\S\S4--6). Everything else is error, and every error
class is named in the ledger (\S10.2).

\section{The two-zone analysis in the dual variable}

This is the paper's central lemma (full derivations in the companion working notes;
conventions anchor-verified to $2\times10^{-14}$).

\begin{lemmafourone}[Parseval]
For real $u_1, u_2 \in L^1(I)$:
\[
\iint_{I\times I} \Phi(\tau-\tau')^2\, u_1(\tau)u_2(\tau')\, d\tau d\tau'
= \int_{\mathbb{R}} g(s)\, F_1(s)\, \conj(F_2(s))\, ds,
\]
$F_j(s) := \int_I u_j(\tau)e^{i\tau s}d\tau$ --- no $2\pi$ factor
in the paperFT convention. In particular the PP block $\mathcal{M}[P_\chi, P_\chi]$
$= \int g(s)|F_\chi(s)|^2\, ds \ge 0$.
\end{lemmafourone}

\begin{lemmafourtwo}[positivity]
$g = \varphi^2\star\varphi^2 \ge 0$ pointwise ([R] \texttt{gv\_nonneg}). Consequently
every $s$-zone $0 \le \int_U g|F|^2 \le \int_{\mathbb{R}} g|F|^2$ separately: the
$\alpha$-zones of the certificate are zones of $s = \alpha\mathcal{L}$, restrictions
of this integral. (What this structure excludes is an $\ell^1$-bound on one half of a
split character sum --- the $(\lVert a\rVert_1 + \sqrt{C}\lVert a\rVert_2)^2$ failure
mode; $n$-range splits with an $\ell^2$/sieve bound on each part are used, and priced,
in Lemma 4.4 and \S5.) (The pointwise positivity --- not the automatic
$\widehat{\psi} \ge 0$ --- is the load-bearing statement; it is forced by
$v = \varphi^2 \ge 0$, which is forced by the Gram realization.)
\end{lemmafourtwo}

\begin{lemmafourthree}[family consumption at constant 1]
With
\[
F_\chi(s) = A_\chi(s) + B_\chi(s), \quad A_\chi(s) = \sum_n a_n(s)\chi(n), \quad
B_\chi(s) = \sum_n b_n(s)\bar{\chi}(n),
\]
$a_n(s) = -(1/2\pi)\Lambda(n)n^{-1/2}D_T(s - \log n)$ (character-independent weights;
no root number, no Gauss sum, no $\log q$), the large sieve applies POINTWISE in $s$
to each squared half; the $\chi/\bar{\chi}$ CROSS term is discharged by
Cauchy--Schwarz in $\chi$ pointwise in $s$ --- the field's own tool for a $\chi(nm)$
bilinear form ([Va, Lemma 6]; [IK, Ch. 7]) --- followed by band separation:
$b_m(s) = \conj(a_m(-s))$ (exactly: $\conj D_T(v) = D_T(-v)$), and every pair
$(n, m)$ is $s$-separated by $\log(nm) \ge \log 4$ BECAUSE $\Lambda(1) = 0$ forces
$n, m \ge 2$. Split the $s$-integral at 0; on $s \ge 0$ the $b$-half carries no
$T$-peak ($|D_T(s + \log m)| \le 2/\log m$), so
$\boldsymbol{\sup_{s\ge0}} \lVert b(s)\rVert_2^2 \le (\log L + M + O(1))/\pi^2$; the
$s < 0$ half is identical under the mirror $s \mapsto -s$, since $g$ is even and
$\lVert b(s)\rVert_2 = \lVert a(-s)\rVert_2$ (on $s < 0$ the roles of the two halves
exchange). The cross term then costs a multiplicative $(1 + \rho)$ with
{\boldmath$\rho \le \sqrt{24/\pi}\cdot\sqrt{(\log L + O(1))/(TL)}$} (the closed form
of the constant 2.7640) --- at design scales
$O(\mathcal{L}^{-(1+r)/2}\sqrt{\log \mathcal{L}})$, a $\delta P \approx 2\times10^{-5}$
at $Q = 10^{25}$ falling to $\approx 3\times10^{-6}$ at $10^{100}$ (this is the
OUT-zone pricing; the in-zone charge is Lemma 4.5's, $\sim$50$\times$ larger and still
sub-minor), well below the smallest budget row. (A factor-$\sqrt{2}$ bookkeeping in
the constant --- whether the mirror doubles the half-integral or reproduces it --- is
left explicitly conservative: even at $\sqrt{48/\pi}$ the term does not register.)
\textbf{Remark (a new obligation, and where it comes from).} CLLR, CIS, Sono and BGST
never meet this cross term --- each arranges it away (a one-sided explicit formula
whose dual prime sum vanishes; $\chi\times\bar{\chi}$ bilinear forms by definition;
$\sigma \leftrightarrow 1-\sigma$ symmetrization). A Gram/Weil realization requires a
REAL density, so both halves are present by construction: the cross term is an
obligation specific to this architecture, and it is discharged by the standard tool
for exactly this object. Since $g \ge 0$, the budget
$X + Q^2 - 1 = Q^2(1 + O(Q^{-\delta}))$ then factors out of the $s$-integral with no
further loss:
\[
\sum_\chi \mathcal{M}[P_\chi, P_\chi] \;\le\; C\cdot|\mathfrak{F}_Q|\cdot(T/\pi)
\cdot\sum_{n\le X} (\Lambda(n)^2/n)\cdot g(\log n)\cdot(1 + o(1)),
\]
i.e.\ family total $\le C \times$ family diagonal --- and the right side matches the
diagonal main term of [R] (their \texttt{prop:PP}), of which ONLY the constant
$(T/\pi)$ transfers (their error term $O(L^2X)$ and their $\lambda \le 1$ hypothesis
do not; the agreement is a consistency check, not an import). No Gallagher lemma, no
hybrid sieve: the hybrid budget $Q^2T + N$ is needed only when $N \gg Q^2$, a regime
this design never enters. Named error terms: the $D_T$-smearing, which is $O(1/T)$
\textbf{in mass-weighted aggregate, zone by zone at $O(\log(TL)/(T\Delta))$ on a zone
of width $\Delta$} (the pointwise version is false near the zone edge and is not
used); and the zone-boundary term, which we state as a lemma:
\end{lemmafourthree}

\begin{lemmafourfour}[zone boundary]
Restricting the $s$-integral to $U = \{|s| \le s_0\}$ does not restrict $n$. Split
$a = a' + a''$ at $n = Y = e^{s_0}$ and apply CS-in-$\chi$ + Lemma 6.1 (never an
$\ell^1$ bound) to the $a''$-part: with $R := \int_U g\lVert a''\rVert_2^2 \le
(\lVert g\rVert_\infty/\pi^2)\cdot\mathcal{L}\cdot\log(T\mathcal{L})\cdot(1+o(1))$
(the $D_T$ tails carry mass $4/y$ beyond distance $y$) and
$P := \int_U g\lVert a'\rVert_2^2 = (T/2\pi)\int_0^{s_0} u\, g(u)\, du\cdot(1+o(1))$,
the relative inflation of the in-zone form is at most
$2C\sqrt{R/P} = O(\sqrt{C \log(T\mathcal{L})/(T\mathcal{L})})$ ---
$\approx 5\times10^{-4}$ at $Q = 10^{100}$ at the design of record (the
$1.5\times10^{-3}$ formerly quoted was the $r = 3$ figure); budget row $L_7$.
$\blacksquare$
\end{lemmafourfour}

\begin{lemmafourfive}[in-zone cross term]
The $\chi/\bar{\chi}$ cross term of the in-zone form --- the third piece of
$\sum_\chi \int_U g|F_\chi|^2 = \sum|A_\chi|^2 + \sum|B_\chi|^2 +
2\re \sum A_\chi \conj(B_\chi)$ --- is NOT covered by Lemma 5.2's orthogonality (the
correction of \S4 withdrew the character-sum route for the cross term), and bounding
it by CS-in-$\chi$ + Lemma 6.1 carries the sieve's $C$ even in-zone. Two bounds,
either sufficient. (i) Conservative, and the one the budget charges: the in-zone
cross term inflates the in-zone form by at most $C\cdot\rho_U$, with $\rho_U$ the
cross-term ratio of this section restricted to $U$ --- $\approx 5\times10^{-4}$ at
$Q = 10^{100}$, a $\delta P \approx 1.6\times10^{-4}$ after the in-zone sensitivity;
charged as minor row $L_{12}$. (The out-zone pricing of the cross term, row
$\sim$3$\times10^{-6}$ at $10^{100}$, does NOT cover this --- the in-zone charge is
$\sim$50$\times$ larger and was previously unpriced.) (ii) Sharper, and the
structural reason the term is harmless: in-zone the $a'/a''$ split of Lemma 4.4
restricts both halves to $n, m \le Y = Q^{1-\delta'}$, where the $\ell^1$ $\times$
divisor route that the correction withdrew GLOBALLY is valid LOCALLY --- the failing
case is a $\lambda > 1$ phenomenon ($n$ reaching $X = Q^\lambda$), and in-zone the
effective bandwidth is below 1: the aggregated character sum of the cross pairs is
bounded by $Q\cdot\tau(\cdot)$ (the divisor average of Lemma 5.3), giving cross
$\ll Q\cdot\lVert a'\rVert_1\lVert b'\rVert_1\cdot\operatorname{polylog} \asymp
Q^{2-\delta'}\cdot\operatorname{polylog}$ against the family diagonal
$\asymp Q^2T\mathcal{L}$ --- negligible at the same zone-edge calibration as Lemma
5.3, with the $a''$ tails absorbed by Lemma 4.4's row. $\blacksquare$ (Corollary 3's
in-zone parity projection consumes Lemmas 4.4--4.5 verbatim: its ``in-zone
coefficient 1'' claim rests on them.)
\end{lemmafourfive}

The in-zone breakpoint is $|s| \le (1 - \delta')\cdot\log Q$, not
$(1 - \delta')\mathcal{L}$ (the difference is budget row $L_1$). (Full proofs in the
companion working notes, which replace the withdrawn divisor-bound treatment of the
cross term.) $\blacksquare$

\section{The character-sum arithmetic}

(Full derivations in the companion working notes; the aggregated identity
anchor-verified to integer precision, 9/9 cases.)

\begin{lemmafiveone}
For $(nm, q) = 1$:
\[
\sum^{*}_{\chi \bmod q} \chi(n)\bar{\chi}(m)
= \sum_{f\mid q,\, f\mid(n-m)} \mu(q/f)\varphi(f)
\]
--- the standard primitive-character
orthogonality relation, in the form displayed in [CIS1, \S3].
\end{lemmafiveone}

\begin{lemmafivetwo}[aggregated, exact]
For $n \ne m$: $\sum_{q\le Q}\sum^{*}_{\chi} \chi(n)\bar{\chi}(m) =
\sum_{f\mid(n-m),\, f\le Q,\, (f,nm)=1} \varphi(f)\cdot M_{nm}(Q/f)$, with
$M_{nm}(y)$ the coprimality-restricted Mertens function. (Moduli with $(q, nm) > 1$
contribute zero since $\chi$ vanishes off coprimality.) Two-line consequence:
$|\sum_{q\le Q}\sum^{*}_{\chi} \chi(n)\bar{\chi}(m)| \le Q\cdot\tau(|n-m|)$, and
$\le Q\cdot\tau(n-1)$ for the linear sums. \textbf{No PNT is used}; $M(y) = o(y)$
sharpens constants only ([CIS1]'s own $\Delta' + \Delta''$ split is this
optimization).
\end{lemmafivetwo}

\begin{lemmafivethree}[zone bound]
$S(Y) := \sum_{n\ne m\le Y} \Lambda(n)\Lambda(m)(nm)^{-1/2}\tau(|n-m|) \le
C\cdot Y(\log Y)^3$ unconditionally (divisor average --- no pointwise
$\tau \le Y^\varepsilon$, whose $\varepsilon$ is NOT $o(\delta')$ and would poison the
zone edge). Hence the low--low zone ($n, m \le Y = Q^{1-\delta'}$) family
off-diagonal is $\ll Q^{-\delta'}\mathcal{L}\cdot(\text{family diagonal})$ --- and
the calibration $\delta' = K\cdot\log \log Q/\log Q$ needs $\boldsymbol{K \ge 2}$: at
$K = 1$ the relative error is $\Theta(1)$, not $o(1)$, and $K = 2$ is the proved
minimum (the companion working notes); the design adopts $K = 3$, a conservative
choice, not a forced one (the budget constant $s(r + K)$ prices each unit of $K$ at
0.5073). The payoff cost $0.5073\cdot\delta'$ sits inside Theorem 1's rate.
\end{lemmafivethree}

\begin{lemmasfivetwoprime}[the $n + m$ variants]
Lemma 5.1 applied at the congruence $-n \equiv m \pmod f$ and aggregated exactly as
in Lemma 5.2 gives the identical identity with the divisor sum over $f \mid (n + m)$;
and since $n, m \ge 2$ forces $n + m \ge 4$ with no excluded diagonal ($n = -m$ is
impossible for positive $n, m$), Lemma 5.3's divisor average applies verbatim with
$\tau(n + m)$ --- the $n + m$ case is, if anything, easier. These are the forms
Corollary 3's in-zone parity projection consumes (\S1.1, \S12.3).
\end{lemmasfivetwoprime}

No case split between the dyadic and full families, and no subfamily sieve, occurs
anywhere in this section --- Lemmas 5.2/5.2$'$ are uniform over $q \le Q$ (the parity
SUBFAMILY appears only in Corollary 3's projector, which consumes them at full-family
level).

\section{The large sieve}

\begin{lemmasixone}[{Gallagher [Ga67]; Montgomery--Vaughan [MV73], [MV74] --- the
additive $N - 1 + \delta^{-1}$ input is [MV74]'s Hilbert inequality, the
multiplicative deduction is classical; see [IK, Thm 7.13] or [Va, Thm 4]; in the
exact form quoted, [PT25, (1.1)]}]
\[
\sum_{q\le Q}\sum^{*}_{\chi \bmod q}\Bigl|\sum_{n\le N}a_n\chi(n)\Bigr|^2
\le (N + Q^2 - 1)\sum_n|a_n|^2.
\]
\end{lemmasixone}

This is the paper's only sieve input, consumed in Lemma 4.3 (pointwise in $s$),
Lemma 8.1 (pointwise in $\tau$), and nowhere else. The multiplicative form is not in
the Lean tree of [R], but its additive engine is: the Montgomery--Vaughan generalized
Hilbert inequality is proved in the artifact (\texttt{Zeta23/MV/}), so the
from-scratch build is the multiplicative deduction (duality + Gauss sums + primitive
decomposition) on top of an in-tree input --- still the largest single build if this
paper's result is to sit beside the record's standard of evidence (\S13).

\begin{remarkenv}[the $q/\varphi(q)$ weight, declined on the record]
The sharp classical form carries a weight:
$\sum_{q\le Q}(q/\varphi(q))\sum^{*}_{\chi}|\cdot|^2 \le (N + Q^2)\sum|a_n|^2$
[Va, Thm 4; IK, Thm 7.13; CIS1, (1.4)]. Dropping $q/\varphi(q) \ge 1$, as Lemma 6.1
does, is valid but not free: under the $(q/\varphi(q))\varphi^{*}$-measure the family
constant would fall to $C_w \approx 4.174$
($\langle q/\varphi(q)\rangle_{\varphi^{*}} = 1.2965$, measured at $Q = 2\times10^6$),
worth $\approx +1.29$ points at $q \le Q$ and $\approx +1.01$ dyadic. We decline it
deliberately: the weighted object is a $q/\varphi(q)$-weighted family average --- a
different statement class from the sharp-count corollaries whose like-for-like
comparisons (+14.99, +8.75) this paper is built to make --- and Lemma 5.2's exact
identity, \S12.2's conductor averages and the shell arithmetic are all stated for the
unweighted count. The option is recorded; the price of the cleaner statement is
$\approx 1.3$ points.
\end{remarkenv}

\section{The tail}

(Full derivations in the companion working notes; anchors: ramp constants
$k = 1\dots8$; envelope on the true $\varrho_2$; the closing ladder; the 60-block
pair-vs-scalar toy.)

\subsection{The Gevrey transform bound}

For any GevreyProfile$(2, A, B)$ taper with $2w \le L$, all $z \in \mathbb{C}$:
$\lVert\widehat{\varphi}(z)\rVert \le e^2\cdot\max(2Bw, L)\cdot e^{|\im z|\cdot L/2}
\cdot\exp(-(2/e)\sqrt{w|z|/A})$ --- from [R]'s proved ramp lemma
$\lVert\varphi^{(k)}\rVert_1 \le 2Bw(A/w)^k k^{2k}$ by $k$-fold parts and
optimization ($k = \lfloor\sqrt{r}/e\rfloor$; the integer floor absorbed by the
small-$r$ branch). The general limitation is the Ingham/Paley--Wiener
logarithmic-integral criterion [In34, PW, H\"o, Thm 1.3.5]: a $Q$-power of decay at
distance $D_0$ with bounded ramp fails for EVERY compactly supported taper; for the
Gevrey-2 class the constraint is $wD_0 \gtrsim L^2$, met by the reference design with
$wD_0 = 3\mathcal{L}^2(\log \mathcal{L})^2$ and by the design of record (\S10.3) with
$wD_0/\mathcal{L}^2 = 8.51$ at $Q = 10^{25}$, 3.80 at $10^{100}$, 2.87 at $10^{300}$,
decreasing to the asymptote $\approx 2.4$ --- the $(1 + o(1))$ is large at headline
scales, and \texttt{table2\_check.py} prints the actual values. (A near-optimal
non-quasianalytic class would relax $L^2$ to $L\cdot\operatorname{polylog}$; we reuse
[R]'s profile and widen the buffer instead. The record's own hard-coded choice
$D_0 = \sqrt{T}$ is not merely suboptimal here: forcing it makes the ramp row
$\Theta(\mathcal{L}^{1-r/2})$, which breaks the theorem's rate class for every
$r < 4$ --- so the $D_0$ generalisation, a field of the \texttt{ParamsQ}
re-parameterisation in the formalization plan, is required, not cosmetic.)

\subsection{The mirrored tail lemma}

Each squared Gram-column entry carries BOTH taper factors at the same argument
([R]'s own $|\gamma - \tau_k|^{-4}$ structure --- the two-factor convention is the
faithful mirror), giving suppression $\exp(-(4/e)\sqrt{w\cdot\operatorname{dist}/A})$
per entry against the off-line amplification $e^{L/2} = X^{1/2}$
($|\im \gamma_\rho| \le 1/2$ against $\widehat{\varphi}$ of exponential type $L/2$
--- the unavoidable unconditional price). Row sums by exponential telescoping, window
sums by the $q$-uniform local count (\S9); the closing condition
$(4/e)\sqrt{wD_0/A} \ge L/2 + \log(\text{prefactor}) + \log(L/\eta)$ holds with
$\Theta(\mathcal{L} \log \mathcal{L})$ to spare at the reference design; the design
of record sits ON the constraint by construction --- the optimiser drives $D_0$ to
the boundary, so the margin is 0 nats at every $Q$ (\texttt{table2\_check.py}), which
is admissible: the condition is an inequality with the $j$-sum prefactor and the
$\log(L/\eta)$ slack already inside it, and equality delivers exactly the target
$\vartheta_0$ (the former ``+3\% at $10^{100}$'' pricing referred to the interior
optimum that the $w \ge 1$ clamp replaced, and is withdrawn). Sensitivity, computed
(\texttt{ext\_checks.py}): REQUIRING a margin of 1 nat moves $P_{\mathrm{eff}}$ by
$\le 0.003$ at $10^{25}$ and $\le 10^{-4}$ from $10^{100}$ on; even a forced margin
of $0.1\cdot L$ ($\approx 28$ nats at $10^{100}$) costs only 0.021/0.002/0.0004 at
$10^{25}$/$10^{100}$/$10^{300}$ --- the zero-margin design is a boundary convenience,
not a knife edge. Either way $\lVert\widetilde{E}_\chi\rVert \le \vartheta_0 \to 0$
faster than any required rate, uniformly over the family. The window $j$-sum is
collected here at crude explicit constants; the honest collection (per-window
decrement $0.350/\mathcal{L}$ at the reference design,
$\sum_j = 2.86\mathcal{L}\cdot(1 + o(1))$; $\approx 0.29/\mathcal{L}$ at the design
of record, same order) is in the companion working notes; the crude form exceeds it
by the factor $0.199\mathcal{L}$ --- conservative, never wrong.

\subsection{Aggregation without a $Q$-power}

Over the direct sum: the operator norm is a MAX (the Weyl consumer pays nothing); the
trace is additive,
$|\tr \widehat{E}_{\mathrm{fam}}| \le |\mathfrak{F}|\vartheta_0/(aL) =: B_{\mathrm{tr}}$;
the Frobenius norm is $\sqrt{\phantom{x}}$-additive,
$\lVert\widehat{E}_{\mathrm{fam}}\rVert_F \le \sqrt{|\mathfrak{F}|}\cdot\vartheta_0/(aL)
=: B_F$. The pair perturbation lemma
($4 \tr \widehat{G} - \lVert\widehat{G}\rVert^2_F - [4B_{\mathrm{tr}} +
2B_F\lVert\widehat{G}\rVert_F + B_F^2] \le 4 \tr \widehat{A} -
\lVert\widehat{A}\rVert^2_F$; three lines, generalizing [R]'s scalar version, which is
the case $B_{\mathrm{tr}} = B_F$) feeds Proposition 3.1's (iv) with per-block
$\vartheta_0 \to 0$ only: \textbf{no negative power of $Q$ is demanded anywhere.}
(The verbatim scalar interface would demand
$\vartheta_0 = o(Q^{-1}\cdot\operatorname{polylog})$; the pair costs nothing and
removes it.)

\section{The ends}

(Full derivations in the companion working notes; anchors on the real 60-character
family.)

\begin{lemmaeightone}[family mean square]
$B_{\mathrm{fam}}(\tau)^2 := \sum_\chi \nu_{X,\chi}(\tau)^2 \le
c_1^2Q^2(\mathcal{L} + \log^{+}(|\tau|/4T) + C_0)^2$, with explicit absolute
$(c_1, C_0)$: the $\mu$-part by digamma envelopes; the $P$-part by Lemma 6.1 at fixed
$\tau$ (the $\tau$-twist leaves $\lVert a\rVert_2$ invariant --- the prime part never
touches the $\log^{+}$ growth). \textbf{No $X$-dependence} --- the standard
large-sieve substitution of a family mean value for a pointwise bound [Mo71, Ch. 12;
IK, Ch. 10; Ga67], applied at the $(B, \nu)$-generic seam of [R]'s ends layer, where
the $T$-aspect's pointwise $B = l + 4\sqrt{X}$ (the $\lambda \le 1$ wall) is replaced
by $B_{\mathrm{fam}} \asymp Q\mathcal{L}$.
\end{lemmaeightone}

\begin{lemmaeighttwo}[bilinear ends bound]
\[
\sum_\chi (\mathcal{E}_1 + \mathcal{E}_2)(\nu_\chi) \le \iint W(\tau, \tau')\,
B_{\mathrm{fam}}(\tau)B_{\mathrm{fam}}(\tau'),
\]
with $W$ the $\nu$-free majorant
kernels of [R] --- quoting the CAP-FREE L-intrinsic constants
($|\mathcal{E}_1| \le 4(90 + 32c\varrho^2)\cdot L^3B^2\cdot(L + l)$;
$|\mathcal{E}_2| \le C_2'\cdot L^3B^2\cdot(1 + \log L)(L + l)$), because the capped
forms' hypothesis $L \le 2l$ --- innocuous in the $T$-aspect, where $L = \lambda l$
--- FAILS here by a factor 10--60. Family relative order:
{\boldmath$\Theta(\mathcal{L} \log \mathcal{L}/T)$} (settling a flagged question from
the companion working notes in favour of the larger figure: the discrepancy was the
capped-vs-cap-free shape, not bookkeeping); budget row
$6\mathcal{L} \log \mathcal{L}/T$, 0.4\% of the budget --- nothing downstream moves.
The hat-normalization bookkeeping reproduces the record's $\lambda \le 1$ ends wall
exactly in the $T$-aspect (consistency anchor). The bilinear write-out with pinned
constants ($c_1 = 0.41$, $C_0 = 6$) is in the companion working notes; its Lean-level
instantiation is deferred to the formalization.
\end{lemmaeighttwo}

\section{The archimedean and explicit-formula layer, uniformly in $q$}

(Full derivations in the companion working notes.) The Weil explicit formula
$Gz(\chi) = Gp(\chi)$ with density $\nu_{X,\chi}$ and no pole term; RvM-$\chi$ with
absolute constant; the local count $N_\chi(t, t+1] \le A_0\cdot\log(q(|t|+3))$ with
ABSOLUTE $A_0$. Every analytic constant in the chain is already explicit-and-linear
in $q$ in [R]'s own development ($\lVert L(s,\chi)\rVert \le q\lVert s\rVert/\re s$;
$\lVert L(w,\chi)\rVert \le 8q(|w.\mathrm{im}|+3)$); the uniformity is re-packaging,
not new analysis. The buffer count
$N_{\mathrm{II},\chi} \le 3A_0D_0\cdot\text{log-scale}$ gives the $L_4$ budget row.
\textbf{The $\lambda \le 1$ / $X \le T$ audit (full file-level table in the companion
working notes):} [R]'s \texttt{Params.Valid} bundles \texttt{lam\_le\_one}, so the
reuse obligation is the \texttt{ParamsQ} weakening (budgeted in the formalization
plan), and \texttt{Tail/} is NOT ``reused verbatim'' (\texttt{theta0\_le} hard-codes
$D_0^2 = T$) --- both corrected from QR.4's first version; the moment core, the ends
\texttt{\_L} variants, the Taper machinery and the uniform local count are cap-free
as audited. The \texttt{LocalCountChi}/\texttt{MainTermChi} re-runs with explicit
constants: the companion working notes (the local count is [R]'s own uniform
theorem); the \texttt{ParamsQ} class is formalization-track.

\section{Assembly, the budget, and the proof of Theorem 1}

\subsection{}
Proposition 3.1's proof is the elementary moment algebra of [R]'s certificate with
the pair split (the companion working notes). \textbf{Remark (deferred to the
formalization):} its Lean-adjacent $\varepsilon$-filter threading.

\subsection{The ledger}
Every explicit-formula error term is classified
[P]/\allowbreak[R]/\allowbreak[CS]/\allowbreak[LS]/\allowbreak[X]/\allowbreak[M]
(19 error
classes in 17 ledger rows; no term with $\tau$-mass $\asymp T$ outside the positive
form); each row states its normalization (coefficient/tilde/hat). The aggregation
discharge: exactly three components are bounded only in family sum (the PP
off-diagonal, the $\mu P$ cross, the ends), each consumed by a certified family
mechanism (\S\S4, 5, 8); no per-block-then-multiply step exists (tabulated in the
companion working notes).

\subsection{The budget (the jointly-optimised design of record)}
Principal terms: zone breakpoint $s\cdot(\log T)/\mathcal{L}$ and zone edge
$s\cdot\delta'$ ($s = 0.5073$, the measured small-$\delta$ payoff SECANT at
$C = \pi^4/18$ --- the secant at $\delta \approx 0.005$; the $\delta \to 0$ slope is
0.5042, and Table 2 uses the larger secant at the actual design offset --- these two
pin the rate class); ramp: the diagonal sandwich $(L-2w)^3/6$ gives relative deficit
$1 - (1 - 2w/L)^3 \approx 6w/L$ in the $\frobSq$ PP-diagonal main (the sandwich is
[R]'s, \texttt{PrimeSideB.lean}:384/:386, cited with the term enumeration in the
companion working notes; the coefficient expansion
$1 - (1 - 2w/L)^3 = 6w/L - 12(w/L)^2 + 8(w/L)^3$ is displayed here, so $6w/L$ is
sharp and charging it in full over-charges for every $w > 0$) --- charged
conservatively in full as $r_2 = 6w/L$, sign restated (the tapered bracket
under-weights the $|\alpha| \le 1$ mass); the previously quoted $4w/L$ ([R]'s
\texttt{calE} shape) and any ``$4r_1$'' reading are retired; buffer $6D_0/T$ ---
priced at the SHARP zero density ($N_{\mathrm{II}}$ at the true $T/2\pi$-per-$\chi$
scale), not at the proved absolute $A_0$ of \S9, which is larger: charging the proved
constant would move the finite-$Q$ thresholds, never the asymptotics, and \S10.4
states which convention its table uses; ends (\S8) --- plus the minor rows (conductor
spread $O(1/\mathcal{L})$; Lemmas 4.4--4.5; the cross term). \textbf{The ramp and
buffer leave the leading order entirely under the balanced choice}: minimising the
ramp-plus-buffer rows ($6w/L + 6D_0/T$) subject to the honest Gevrey closing
condition $(4/e)\sqrt{wD_0/A} \ge L/2 + \log(\text{prefactor}) + \log(L/\eta)$ gives
$w^* = \Theta(\mathcal{L}^{(3-r)/2})$ --- $\Theta(1)$ at $r = 3$, slowly decreasing
for $r > 3$ --- CLAMPED at the regime floor $w \ge 1$ carried by the reused lemmas
([R]'s \texttt{Params.Valid} field \texttt{one\_le\_w}; side conditions in the
companion working notes), which binds at the design of record from
$Q \approx 10^{50}$; so the design sets $w = \max(1, w^*)$ ($= 1$ at $Q = 10^{100}$;
the previously quoted $\approx 2.3$ was the $r = 3$ optimum under the retired $4w/L$
objective, and is withdrawn), $D_0$ balanced accordingly, every side condition
satisfied (the admissibility condition $8w \le L$, the QT.a optimisation, $wD_0 \gg e^2A$), and ramp + buffer
$= O(1/\mathcal{L})$ --- the clamp contributes $6/(\lambda\mathcal{L})$, the same
order with constant $6/\lambda^*$. One accounting question is stated openly for
review: [R]'s aggregate \texttt{calE} carries a trace-side $w/L$, which is NOT a
separate charge in our accounting --- the hat normalisation $a := L^{-1}\int\varphi^2$
is computed for the TAPERED $\varphi$, so the taper's first-order mass deficit
cancels in $\tr \widehat{G}$ against $a\mathcal{N}$, and the trace-side residue is
the buffer/tail rows already charged; if a reviewer rules otherwise, the fully
conservative charge is $r_2 + 4r_1 = 10w/L$, costing $\le 0.012$ at $Q \ge 10^{100}$
with the rate class and constant untouched. Total $= O(\log \log Q/\log Q)$, pinned
by the two zone terms. \textbf{Remark (the constant, computed, not proved).} The
asymptotic budget constant is $s\cdot(r + K)$ at $T = (\log Q)^{r+\varepsilon}$ and
$\delta' = K \log \log Q/\log Q$: \textbf{\boldmath$\approx 3.3$ at the design
$(r, K) = (3.5, 3)$} (2.54 at $(r, K) = (3, 2)$, $K = 2$ being the minimum \S5
actually proves; $\approx 4.6$ at $(6, 3)$, the choice that additionally sits inside
[CIS2]'s window floor $(\log Q)^{6}$) --- a computed property of the
jointly-optimised construction, NOT in the theorem, whose implied constant is
effectively computable. (With the $w \ge 1$ clamp the measured constant at
$\log Q = 2.3\times10^6$ is 3.56, converging to $s(r + K)$ from above: the clamp's
$6/\lambda\mathcal{L}$ sits a $\log \mathcal{L}$ below the leading term.)
\textbf{Distinct objects:} this ramp row (the taper's diagonal-mass cost, first order
in $w/L$) and \S11's ramp cost (the variational profile's edge perturbation, third
order --- provably $\Theta(\rho^3)$) are different quantities; both appear because
the taper acts once on the window and once at the free boundary.

\subsection{Table 2 --- for orientation only}
[the only effective statements in the paper; they appear in this remark and nowhere
else --- not in the abstract, not in the introduction]: Design $(r, K) = (3.5, 3)$
--- $T = (\log Q)^{3.5}$, $\delta' = 3 \log \log Q/\log Q$; zone cost from the
measured $P(\delta)$ secant at the design offset; the ramp row charged at \S10.3's
conservative $r_2 = 6w/L$ with the $w \ge 1$ regime floor enforced; the buffer row at
the sharp zero density (\S10.3 --- the one row NOT charged at its proved constant);
ALL rows carried including the minor ones, Lemma 4.5's in-zone cross row now among
them (computed by the shipped \texttt{table2\_check.py}; \texttt{log\_table2.txt}):

\begin{center}
\bfseries\boldmath
\begin{tabular}{lccccc}
$Q$ & $10^{25}$ & $10^{100}$ & $10^{300}$ & $10^{1000}$ & \\
$P_{\mathrm{eff}}$ & $+0.20$ & $+0.61$ & $+0.68$ & $+0.71$ & $\to 0.72128$ \\
\end{tabular}
\end{center}

\noindent\textbf{\boldmath non-vacuous from $\approx 10^{20}$; $> 0.5$ from
$\approx 10^{51}$; exceeding the per-character record from $\approx 10^{236}$.} Under the sharper $4w/L$ reading every figure
improves --- by +0.044 at $10^{25}$ and by $\le +0.007$ from $10^{100}$ on --- and
the thresholds fall to $\approx 10^{18}$ / $10^{48}$ / $10^{220}$. Thresholds below
$\approx 10^{30}$ are extrapolations of an asymptotic model and should be quoted as
such, if at all. We have not attempted to optimize the implied constant; an effective
version of Theorem 1 would require explicit forms of the inputs of \S\S5--9 and is a
separate undertaking. (For calibration: no prior FAMILY-AVERAGE result in this line
--- CIS, CLLR, Sono --- states any effective range; the computed threshold here is
believed to be, if anything, smaller than those constructions would yield.)

\subsection{}
Combining \S\S4--9 into Proposition 3.1's hypotheses at $\lambda = \lambda^*$ proves
Theorem 1 at each fixed $\varepsilon > 0$, with the rate carried by the budget rows
$r_i$ --- no $\varepsilon$-limit is taken, and $Q_0(\varepsilon)$ is the largest of
\S\S5--9's thresholds; the dyadic corollary runs the same argument at
$C = 2\pi^4/27$ (its finite-$Q$ arithmetic from the measured dyadic $P(\delta)$
secant at the design offset --- 0.64 at the $10^{100}$ offset; the small-$\delta$
secant 0.5485 is not usable at finite $Q$).

\section{The variational layer}

The framework is Sono's [So16] (the RKHS/ODE class; sign-dual:
Chirre--Gon\c{c}alves--de Laat); the $C$-penalised free-boundary solution is a
calculation within it: minimize
$B(v) = \psi(0) + \int_{|\alpha|\le1}|\alpha|\psi + C\int_{1<|\alpha|\le\lambda}
|\alpha|\psi$, $\psi = v\star v$, over $v \ge 0$, $\int v = 1$,
$\supp \subseteq [-\lambda/2, \lambda/2]$; $P = 2 - \min B$. Optimal $v$:
$\cos(\sqrt{2}t)$ in the bulk; edge zones from the quartic
$(\rho^2 + 2)^2 = (C-1)^2(1 - \rho^2)$ with the reflection collapse; free boundary
$v(b^*) = 0$. Values (ten digits; three calibration gates: MT $= 0.672500703679412$;
flat-$v$ $4/3$; and [So16]'s PUBLISHED GRH benchmark reproduced on the strictly
better kernel $F = \min(|\alpha|, 1)$ at $\lambda = 2$ --- 0.9322826239 against
Sono's 0.93228262, \texttt{gate3\_so16.py} --- a gate that lands on someone else's
published number. One honest caveat, stated because a referee will notice it: all
three gates are special cases --- fixed kernel or a $C$-limit --- so they certify the
SOLVER against outside ground truth; the finite-$C$ free-boundary regime the headline
uses is checked separately, by the closed-form/Nystr\"om agreement to twelve digits
and the strict KKT complementarity of \texttt{payoff\_hp2.py}):
$\pi^4/18 \to 0.7212835668$ at $\lambda^* = 1.2507321515$;
$2\pi^4/27 \to 0.7099167448$ at $\lambda^* = 1.1931581210$; gain positive for EVERY
finite $C$ (no break-even); $\lambda = 1^{+}$ slope $\approx 0.68$ (weakly
$C$-dependent: 0.6847 at $C = 1$, 0.6819 at $C = 20$). Remark (an INADMISSIBLE
limit, for orientation only): at $C = 1$ --- unattainable, since $C \ge \pi^4/18$ ---
the plateau is exactly $2 - \pi/(2\sqrt{2}) = 0.8892792655$ at
$\lambda = \pi/\sqrt{2} = 2.2214$, which also exceeds the sieve's own range
$\lambda < 2$; the constrained $\lambda = 2$ ceiling is 0.88836541. Neither is
comparable to [So16]'s GRH 93.22\% ($M = 0.93228262$), which optimises the strictly
better kernel $F = \min(|\alpha|, 1)$ --- the third calibration gate above. Two
documented solver traps (the kernel jump at $|\alpha| = 1$; the constrained-window
junction) --- any new implementation must pass the gates. The taper's ramp does not
disturb the optimum, and provably: $v^*$ vanishes linearly at the free boundary, so a
ramp of relative width $\rho$ gives $\lVert\delta v\rVert_1 = O(\rho^2)$ and
$\lVert\delta v\rVert_2^2 = O(\rho^3)$; $B$ is stationary at $v^*$ for
mass-preserving perturbations, and of the second variation the kernel part pairs two
$\ell^1$-masses ($O(\rho^4)$) while $\psi(0) = \int v^2$ is an $L^2$-norm --- so
\textbf{\boldmath$\Delta P = \Theta(\rho^3)$, with the $\int v^2$ term as the
reason} (stable across three ramp shapes and three grids; measured exponent
3.03--3.09). The free-boundary condition itself protects the profile.

\section{Normalisation of the $q \le Q$ family}

\subsection{}
The consumed object is a large-sieve majorant of the family second moment of a
prime-supported Dirichlet polynomial at the common scale $\mathcal{L}$ --- defined by
the construction, single-scale by fiat of the common taper (itself forced by Lemma
6.1's single coefficient vector). It is not [CLLR]'s $F_\Phi$.

\subsection{}
The per-character scales that DO survive ($\mu_q$, the trace, the denominators
$N_\chi$) enter as bounded powers of $c_q = \log q/\mathcal{L}$ under
$\varphi^{*}$-weights, and the certificate consumes a ratio of sums, never an average
of ratios: $\langle c^2\rangle/\langle c\rangle$ and $1/\langle c\rangle$ are
$1 + O(\log \log Q/\log Q)$ (driven by $\log T/\mathcal{L}$; the conductor-spread
contribution alone is $O(1/\mathcal{L})$ --- the same order as the rate, and absorbed
by it). The conductor averages, CHARACTER-COUNT-weighted ($\varphi^{*}$-weighted), by
partial summation against $\sum_{q\le x}\varphi^{*}(q) = (18/\pi^4)x^2 + O(x \log x)$:
$\langle \log q\rangle = \log Q - 1/2$ over $q \le Q$, and
$\log Q + (\log 2)/3 - 1/2 = \log Q - 0.26895$ over $(Q/2, Q]$ --- for contrast, the
UNWEIGHTED averages are $\log Q - 1$ and $\log Q + \log 2 - 1 = \log Q - 0.30685$;
the latter is the quantity Fiorilli--Miller compute, cited here as the precedent for
computing (rather than discarding) the conductor-averaging correction, with
[Mi, Rem. 3.2] for the family-wide-scale recommendation ([Mi, (2.8)] itself is stated
for subsets of $[N, 2N]$ and does not cover $q \le Q$ --- the two-line derivation
above is what does). Every subfamily this paper states pays its own constant,
itemised: Corollary 2 (dyadic: $C = 2\pi^4/27$, conductor average
$\log Q - 0.26895$) and Corollary 3 (even dyadic: $C = 4\pi^4/27$ via sieve
positivity, the parity restriction being exactly a subfamily localisation ---
\S12.3).

\subsection{}
The Sono comparison is Corollary 3's, and the parity restriction is priced, not
waved: the passage is the two-zone hybrid of Corollary 3 --- parity projection
IN-ZONE ($\tfrac{1}{2}\sum^{*}(1 + \chi(-1))$, the $n \equiv -m$ orthogonality by
Lemma 5.2$'$, valid there because $n + m \le 2Q^{1-\delta'} \ll Q$), positivity
OUT-ZONE at $C \to 2C$ (the projection route fails there at $\lambda > 1$ by
$Q^{\lambda-1}$ --- the identical mechanism as the \S4 cross term, with no band
separation available since both halves share a band --- and no parity-restricted
sieve at budget $\tfrac{1}{2}(N + Q^2 - 1)$ appears in the literature). The mechanism
generalises exactly as far as its TWO obligations do: a subfamily of positive
relative density $\alpha$ gets the argument at out-zone constant $C/\alpha$ --- above
the $\lambda = 1$ floor 0.672500703679, since \S11's gain is positive for every
finite $C$ --- only if (i) its conductor distribution matches the full family's to
leading order (else \S12.2 must be redone, as Corollary 2 does with its own $C$ and
its own $\langle \log q\rangle$), AND (ii) it admits an in-zone orthogonality
projector expressing its sum through full-family instances of Lemmas 5.2/5.2$'$, as
the parity classes do. A general density-$\alpha$ subfamily has no such projector,
and positivity in-zone gives coefficient $1/\alpha$ --- the 0.3693 trap named in
Corollary 3: selecting within each modulus the half of the characters with the
largest in-zone mass satisfies (i) exactly and fails the floor, so (ii) is not
removable. The proved cases are exactly three: even and odd (the parity projectors),
and dyadic (conductor-restricted, with \S12.2 redone). Residual mismatches with
[So21]: the smooth weight (inflates $C$ by $3/(2\int Wx \, dx)$ if mirrored ---
stated, not paid, since our sharp cutoff is the stronger statement) and nothing else:
$[T, 2T]$ matches [So21] and [R] exactly. ([CIS2]'s window is $|\gamma| \le T$ with
floor $(\log Q)^{6}$ --- our $T = (\log Q)^{3+\varepsilon}$ sits BELOW that floor; at
$r \ge 6$ the theorem enters [CIS2]'s window class too, at budget constant
$\approx 4.6$.)

\section{Formalization outlook}

A Lean 4 formalization is planned and scoped: most of \S\S3--10 mirrors [R]'s tree at
file-level pointers recorded in the companion working notes; the from-scratch builds
are Lemma 6.1 (the sharp multiplicative large sieve) and the re-parameterisation of
[R]'s parameter structure ($\lambda > 1$, free $D_0$). It is not claimed as done;
when complete it will be held to [R]'s own standard (sorry-free, axiom-audited).

\section*{Disclosure of generative-AI use}

This work was produced substantially by a generative AI system --- Claude (Anthropic)
--- working under the direction of the named author: the derivation and optimisation
of the construction of \S\S4--12, the numerical work of \S\S10--11 and the ancillary
scripts, and the drafting of the exposition are all substantially AI-generated, and
this is disclosed in accordance with arXiv's policy on generative AI language tools.
No AI system is listed as an author. Verification of the mathematical content has
been by process rather than by the author's own line-by-line reading: repeated
rounds of independent adversarial review (AI-conducted, on blinded bundles,
including fully independent sessions with no project history), numerical
calibration gates that reproduce two published third-party constants (\S11),
reproducible ancillary scripts for every quoted figure, and the planned Lean
formalization of \S13, which is the intended final arbiter. The named author takes
full responsibility for the entire contents of the paper, irrespective of how any
part of it was generated. The preprint [R], from which the six hypotheses H1--H6 of
\S2.1 are imported, is itself an AI-generated work; \S2.1 isolates every statement
taken from it, and its public Lean 4 artifact is the intended means of checking those
statements independently of this paper.

\begin{thebibliography}{BGST2}

\bibitem[BGST]{BGST} S. A. C. Baluyot, D. A. Goldston, A. I. Suriajaya,
C. L. Turnage-Butterbaugh, \emph{An unconditional Montgomery theorem for pair
correlation of zeros of the Riemann zeta-function}, Acta Arith. 214 (2024), 357--376;
arXiv:2306.04799

\bibitem[BGST2]{BGST2} iid., \emph{Pair correlation of zeros of the Riemann zeta
function I: proportions of simple zeros and critical zeros}, arXiv:2501.14545
(unpublished; its Thm 2 carries 0.672500703679 conditionally, under a
narrow-vertical-box hypothesis --- [R]'s contribution is removing it)

\bibitem[BP]{BP} S. Baluyot, K. Pratt, \emph{Dirichlet L-functions of quadratic
characters of prime conductor at the central point}, J. Eur. Math. Soc. 24 (2022),
369--460

\bibitem[CGdL]{CGdL} A. Chirre, F. Gon\c{c}alves, D. de Laat, \emph{Pair correlation
estimates for the zeros of the zeta function via semidefinite programming}, Adv.
Math. 361 (2020), art. 106926

\bibitem[CIS1]{CIS1} J. B. Conrey, H. Iwaniec, K. Soundararajan, \emph{The asymptotic
large sieve}, arXiv:1105.1176 (2011; unpublished --- so listed on Conrey's own
publication page; its (1.4) carries the $q/\varphi(q)$ weight, and its introduction
the convention sentence quoted in \S1.4)

\bibitem[CIS2]{CIS2} iid., \emph{Critical zeros of Dirichlet L-functions}, J. reine
angew. Math. 681 (2013), 175--198; arXiv:1105.1177

\bibitem[CIS3]{CIS3} iid., \emph{The mean square of the product of a Dirichlet
L-function and a Dirichlet polynomial}, Funct. Approx. Comment. Math. 61 (2019),
147--177; arXiv:1808.02879 (published; the mean-value engine under [So21])

\bibitem[CL]{CL} V. Chandee, Y. Lee, \emph{n-level density of the low lying zeros of
primitive Dirichlet L-functions}, Adv. Math. 369 (2020), art. 107185;
arXiv:1706.02848

\bibitem[CLLR]{CLLR} V. Chandee, Y. Lee, S.-C. Liu, M. Radziwi\l\l, \emph{Simple
zeros of primitive Dirichlet L-functions and the asymptotic large sieve}, Q. J. Math.
65 (2014), 63--87; arXiv:1211.6725

\bibitem[Di24]{Di24} M. Dickinson, \emph{Zeros of Dirichlet L-functions near the
critical line}, Mathematika (2024), doi:10.1112/mtk.12239; arXiv:2211.06264 (38.2\%
on the line for the primitive family of a single fixed modulus,
$T \gg q^\varepsilon$)

\bibitem[DPR]{DPR} S. Drappeau, K. Pratt, M. Radziwi\l\l, \emph{One-level density
estimates for Dirichlet L-functions with extended support}, Algebra \& Number Theory
17 (2023), 805--829

\bibitem[FM]{FM} D. Fiorilli, S. J. Miller, \emph{Surpassing the Ratios Conjecture in
the 1-level density of Dirichlet L-functions}, Algebra \& Number Theory 9 (2015),
13--52; arXiv:1111.3896

\bibitem[Ga67]{Ga67} P. X. Gallagher, \emph{The large sieve}, Mathematika 14 (1967),
14--20

\bibitem[Go]{Go} D. A. Goldston, \emph{Notes on pair correlation of zeros and prime
numbers}, in \emph{Recent Perspectives in Random Matrix Theory and Number Theory},
LMS Lecture Note Ser. 322, Cambridge Univ. Press (2005), 79--110; arXiv:math/0412313

\bibitem[GS26]{GS26} D. A. Goldston, A. I. Suriajaya, \emph{Zeta Zeros on the
Critical Line}, arXiv:2511.20059 (expository, v2 5 Feb 2026; its 41.7\%/40.7\% are
[PRZZ]'s constants, for $\zeta$)

\bibitem[GZ]{GZ} P. Gao, L. Zhao, \emph{One level density of low-lying zeros of
families of L-functions}, Compositio Math. 147 (2011), 1--18

\bibitem[H\"o]{Ho} L. H\"ormander, \emph{The Analysis of Linear Partial Differential
Operators I}, Grundlehren 256, Springer (1983), Thm 1.3.5

\bibitem[IK]{IK} H. Iwaniec, E. Kowalski, \emph{Analytic Number Theory}, AMS Colloq.
53 (2004); Thm 7.13 (the multiplicative large sieve), Ch. 7, Ch. 10

\bibitem[ILS]{ILS} H. Iwaniec, W. Luo, P. Sarnak, \emph{Low lying zeros of families
of L-functions}, Publ. Math. IH\'ES 91 (2000), 55--131

\bibitem[In34]{In34} A. E. Ingham, \emph{A note on Fourier transforms}, J. London
Math. Soc. 9 (1934), 29--32

\bibitem[LM]{LM} J. Levinson, S. J. Miller, \emph{The n-level densities of low-lying
zeros of quadratic Dirichlet L-functions}, Acta Arith. 161 (2013), 145--182;
arXiv:1208.0930

\bibitem[Mi]{Mi} S. J. Miller, \emph{Density functions for families of Dirichlet
characters} (unpublished note, 1999)

\bibitem[Mo71]{Mo71} H. L. Montgomery, \emph{Topics in Multiplicative Number Theory},
LNM 227 (1971)

\bibitem[MV73]{MV73} H. L. Montgomery, R. C. Vaughan, \emph{The large sieve},
Mathematika 20 (1973), 119--134

\bibitem[MV74]{MV74} iid., \emph{Hilbert's inequality}, J. London Math. Soc. (2) 8
(1974), 73--82

\bibitem[\"Oz]{Oz} A. E. \"Ozl\"uk, \emph{On the q-analogue of the pair correlation
conjecture}, J. Number Theory 59 (1996), 319--351

\bibitem[PRZZ]{PRZZ} K. Pratt, N. Robles, A. Zaharescu, D. Zeindler, \emph{More than
five-twelfths of the zeros of $\zeta$ are on the critical line}, Res. Math. Sci. 7
(2020), art. 20; arXiv:1802.10521

\bibitem[PT25]{PT25} A. Pascadi, J. Thorner, \emph{Large sieves for $GL_n$ and
applications}, arXiv:2508.14888 (2025; cited for its (1.1), the sharp multiplicative
large sieve)

\bibitem[PW]{PW} R. Paley, N. Wiener, \emph{Fourier Transforms in the Complex
Domain}, AMS Colloq. 19 (1934), Thm XII

\bibitem[R]{R} ``Claude'' (Anthropic), \emph{More Than Two Thirds of the Zeros of the
Riemann Zeta Function Are Simple and On the Critical Line}, preprint, dated 11 August
2026 (released 10 August 2026; revised in place 13 August 2026, per the publisher's
changelog; not posted to arXiv; no journal submission, refereeing, DOI or erratum
known to us); current PDF:
\url{www-cdn.anthropic.com/95c246936988e43127bc6b2ceb7077c1dad2d68e.pdf} (the
10 August version, titled ``\dots Lie on the Critical Line'' and lettered A--E,
remains live at \dots\path{/564f962e60643842f5fcb4a17c9dbc8f608f1c37.pdf}); Lean 4
artifact: \url{github.com/anthropics/zeta-23-lean} (Lean v4.33.0-rc2; Mathlib
51e6992e; Apache-2.0; sorry-free; \texttt{\#print axioms} reports only
\texttt{propext}, \texttt{Classical.choice}, \texttt{Quot.sound}; the artifact
retains the A--E lettering, in which the Dirichlet case is Theorem E)

\bibitem[So16]{So16} K. Sono, \emph{A Note on Simple Zeros of Primitive Dirichlet
L-Functions}, Bull. Aust. Math. Soc. 93 (2016), 19--30 (the GRH low-lying benchmark,
$M = 0.93228262$)

\bibitem[So21]{So21} K. Sono, \emph{Zeros of Dirichlet L-functions on the critical
line}, J. Number Theory (2025; PII S0022314X25000319); arXiv:2105.07422 (2021). The
arXiv page carries no journal reference, which misled earlier drafts of this paper;
the even-primitive restriction is in its (1.8); window floor $(\log Q)^{2}$

\bibitem[Va]{Va} R. C. Vaughan, \emph{The Large Sieve} (Penn State lecture notes,
\path{personal.science.psu.edu/rcv4/LargeSieve.pdf}), Lemma 6, Thm 4

\bibitem[Wu16]{Wu16} X. Wu, \emph{Distinct Zeros and Simple Zeros for the Family of
Dirichlet L-Functions}, Q. J. Math. 67 (2016), 757--779; arXiv:1206.1679
(unconditionally 60.261\% simple --- critical line NOT required --- and 66.433\%
under GRH, over the $\Psi(q/Q)/\varphi(q)$-weighted primitive family; cites [CIS1]
by arXiv number as unpublished)

\bibitem[Wu19]{Wu19} X. Wu, \emph{The twisted mean square and critical zeros of
Dirichlet L-functions}, Math. Z. 293 (2019), 825--865; arXiv:1802.09704 (per
character, $\log q = o(\log T)$: 41.72\% on the line, 40.74\% simple-and-on-line ---
the strongest PUBLISHED per-character unconditional proportions).

\end{thebibliography}

\end{document}
